% ======================================================================
\chapter{Introdução}
\label{cap:introducao}
% ======================================================================
\instrucao{De duas a quatro páginas. Apresente o fenômeno, por que medir a grandeza escolhida e o que já se sabe sobre ela. Termine com o objetivo geral, os objetivos específicos e a organização do texto.}

A velocidade do som em um gás depende da temperatura e da composição do meio, e sua medida é um dos experimentos clássicos de ondas mecânicas \cite{halliday2016}. Em um tubo fechado em uma extremidade, as ondas sonoras formam ondas estacionárias apenas para comprimentos de coluna bem definidos, o que permite obter o comprimento de onda com uma trena e um gerador de frequência \cite{nussenzveig2014}.

\section{Objetivos}

O objetivo geral deste trabalho é determinar a velocidade do som no ar, com a incerteza expressa segundo o GUM \cite{jcgm2008}. Os objetivos específicos são:
\begin{alineas}
  \item montar um tubo de ressonância com coluna de ar ajustável;
  \item medir as posições de ressonância para uma frequência fixa;
  \item ajustar a relação linear entre a posição e a ordem da ressonância;
  \item comparar o resultado com o valor previsto para a temperatura ambiente.
\end{alineas}

\section{Organização do trabalho}

O \autoref{cap:fundamentacao} apresenta o modelo de ondas estacionárias em tubos. O \autoref{cap:metodologia} descreve a montagem e o procedimento. O \autoref{cap:resultados} traz as medidas, o ajuste e a discussão, e o \autoref{cap:conclusao} resume os resultados.
